% =====================================================================
% 统计学论文 LaTeX 骨架（JASA/统计顶刊风格）
% 用途：AI 辅助写作时的「格式锚点」——所有生成内容按此结构填充
% 配套：appendix.tex / supplement.tex / refs.bib
% 编译：xelatex main.tex && bibtex main && xelatex main.tex && xelatex main.tex
% =====================================================================
\documentclass[11pt]{article}

% ---------- 宏包 ----------
\usepackage{amsmath,amssymb,amsthm,mathtools}
\usepackage{graphicx}
\usepackage{booktabs}          % 三线表（统计论文表格标准）
\usepackage{multirow}
\usepackage[margin=1in]{geometry}
\usepackage{setspace}
\usepackage[colorlinks=true,linkcolor=blue,citecolor=blue,urlcolor=blue]{hyperref}
\usepackage[capitalize,noabbrev]{cleveref}   % \cref 自动引用定理/图表
\usepackage{enumitem}
\usepackage{natbib}            % author-year 引用（统计期刊主流）

% ---------- 定理环境（统计论文标准配置） ----------
\newtheorem{theorem}{Theorem}
\newtheorem{lemma}[theorem]{Lemma}
\newtheorem{proposition}[theorem]{Proposition}
\newtheorem{corollary}[theorem]{Corollary}
\newtheorem{definition}[theorem]{Definition}
\newtheorem{assumption}{Assumption}          % 单独编号，按 Assumption 1 引用
\newtheorem{remark}[theorem]{Remark}
\newtheorem{example}[theorem]{Example}
\newproof{pf}{Proof}                          % 或直接用 amsthm 的 proof 环境
\newproof{pf sketch}{Proof Sketch}            % 正文只放证明思路，全文在附录

% ---------- 记号表（AI 生成记号后必须人工核对一次） ----------
\newcommand{\X}{\bm{X}}
\newcommand{\y}{\bm{y}}
\newcommand{\betaHat}{\hat{\bm{\beta}}}
\newcommand{\norm}[1]{\left\lVert #1 \right\rVert}
\newcommand{\E}{\mathbb{E}}
\newcommand{\Prob}{\mathbb{P}}
\newcommand{\ind}[1]{\mathbb{1}\{#1\}}
\newcommand{\dconv}{\xrightarrow{d}}          % 依分布收敛
\newcommand{\pconv}{\xrightarrow{p}}          % 依概率收敛
\newcommand{\aconv}{\xrightarrow{a.s.}}       % 几乎必然收敛

% ---------- 自检标记宏（定稿前必须全部清零） ----------
\newcommand{\TODO}[1]{\textcolor{red}{[TODO: #1]}}            % 待补内容
\newcommand{\AICHECK}[1]{\textcolor{orange}{[AI-待核验: #1]}} % AI生成待人工核验处
\newcommand{\REFCHECK}[1]{\textcolor{orange}{[引用核验: #1]}} % 待核验引用

\title{标题：方法名 + 解决的统计问题 + 方法类型（如：A Penalized Likelihood Approach for \ldots）}
\author{作者\thanks{通讯作者。邮箱。本文代码与数据复现包见 \url{https://github.com/...}。AI 辅助声明见附录 \ref{app:ai-disclosure}。}}
\date{}

\begin{document}
\maketitle

% =====================================================================
% 摘要：150-250 词。结构 = 1句问题背景 + 1句现有方法缺陷 + 1句本文方法
%       + 1-2句理论保证（明确写出收敛阶/分布结果）+ 1句实验结论 + 1句意义
% =====================================================================
\begin{abstract}
TODO：在什么科学问题与数据形态下（背景），现有方法因什么统计性质缺陷而不适用（缺口），
本文提出什么方法（机制一句话），在什么条件下建立了什么理论保证（相合性、渐近正态、收敛速率、推断有效性），
模拟与真实数据实验表明什么（相对基准的量化提升），对什么领域实践有何意义。
\bigskip

\noindent\textbf{Keywords:} 3-6 个关键词，从 \textit{Statistical Methods} 主题词表中选取（如 high-dimensional inference; penalized regression; robust estimation）
\end{abstract}

\newpage
\tableofcontents

% =====================================================================
\section{Introduction}\label{sec:intro}
% JASA 风格 Introduction 的五段式骨架（详见 prompts/P17）：
% (1) 科学问题与数据背景 →(2) 现有方法与统计性质缺口 →(3) 本文方法与机制直觉
% →(4) 理论贡献与主要结果（列点）→(5) 文献脉络 + 全文结构
\subsection{A motivating problem}\label{sec:motivate}
% 用一个具体科学问题/数据集开篇。必须出现：真实数据形态、为什么常规方法失效。

\subsection{Our contributions}\label{sec:contrib}
\begin{enumerate}[label=(\roman*)]
  \item \textbf{方法贡献：} TODO（一句话：提出 X，其机制是 Y）
  \item \textbf{理论贡献：} TODO（在 Assumption \ref{asm:1} 下建立
        $\norm{\betaHat-\beta_0}=O_p(\sqrt{s\log p/n})$ 与 $\sqrt{n}$-渐近正态，支撑后验/置信推断）
  \item \textbf{计算贡献：} TODO（算法复杂度 $O(n p^2)$ 每迭代，R/Python 包开源）
\end{enumerate}
% 注意：贡献列表中每一条都要能在正文/附录中找到对应定理编号或图表编号，见 scripts/notation_check.py

\subsection{Related work}\label{sec:related}
% 按方法脉络分 2-3 小段；每段结尾必须指出与本文的区别（差分句式）。
\REFCHECK{本节所有引用逐条通过 scripts/citation_check.py 三层核验}

\subsection{Organization}\label{sec:org}
% 全文结构 + 附录/补充材料说明。

% =====================================================================
\section{Methodology}\label{sec:method}
\subsection{Problem setup and notations}\label{sec:setup}
% 明确：观测 $(\X,\y)$ 的生成模型；p, n, s 的含义；
% 所有记号在本文第一次出现处定义，全文不得复用歧义。

\subsection{The proposed estimator}\label{sec:estimator}
\begin{equation}\label{eq:obj}
  \betaHat = \arg\min_{\beta\in\mathbb{R}^p}
  \left\{ \frac{1}{n}\sum_{i=1}^n \ell(y_i; x_i^\top\beta)
  + \mathcal{P}_\lambda(\beta) \right\},
\end{equation}
% 公式后必须紧跟：目标函数每一项的统计含义；$\lambda$ 的作用；
% 与既有目标函数的形式差异（差分句式）。

\subsection{Algorithm}\label{sec:algorithm}
% 算法伪代码 + 收敛性说明 + 计算复杂度。
% 伪代码用 algorithm2e 或 algorithmic；每步可执行（允许他人直接复现）。

% =====================================================================
\section{Theoretical results}\label{sec:theory}
% 统计论文理论节的呈现顺序（铁律）：
% Assumptions（编号、可讨论）→ Main theorem → 讨论 Remark（对比文献/速率最优性）
% → 证明策略一句话（证明细节全部放附录）。

\begin{assumption}\label{asm:1}
TODO：设计矩阵的次高斯性 / 稀疏性 $s\log p = o(n)$ / 损失函数的凸性与光滑参数。
每条假设必须有 remark 讨论：是否可放松、与文献假设的强弱比较。
\end{assumption}

% 定理引用假设的规范写法：多条假设时用范围引用（需先补足对应假设环境）
\begin{theorem}\label{thm:main}
在 Assumption \ref{asm:1} 下，若 $\lambda \asymp \sqrt{\log p/n}$，则以概率至少 $1-\delta$，
\begin{equation}\label{eq:rate}
  \norm{\betaHat-\beta_0} \le C \sqrt{\frac{s\log p}{n}}.
\end{equation}
\end{theorem}

\begin{remark}
TODO：速率与 minimax 下界的对齐说明（最优性）；$\lambda$ 选取的实践建议。
\end{remark}

\begin{pf sketch}
TODO：三句话给出证明路线（如：基不等式 + 约束误差 + 峰度引理），细节见附录 \ref{app:proof}。
\end{pf sketch}
% 证明逐行审查用 skills/stat-theory-check，核对清单见 checklists/C2-理论检查清单.md

% =====================================================================
\section{Simulation studies}\label{sec:simulation}
% 模拟实验四要素（缺一不可，详见 checklists/C3）：
% 设计网格 (scenarios) | 重复次数与随机种子 | 评价指标 | 呈现方式（表/图）
\subsection{Simulation design}\label{sec:simdesign}
% 场景网格：n ∈ {200,500}, p ∈ {100,1000}, 误差分布 ∈ {N, t_3, ...}，
% 重复 500 次；随机种子固定并公开；数据生成过程 DGP 逐条列出。

\subsection{Results}\label{sec:simresults}
% 表格遵循 three-digits rule（见 docs/08）；报告均值±标准差；
% 每张表/图必须在正文中被文字解读，且解读与数字一致（scripts/fig_table_check）。

% =====================================================================
\section{Real data analysis}\label{sec:real}
% 数据来源与可得性声明 → 预处理（可复现）→ 方法应用 → 结果解读 → 与领域文献结论对比。

% =====================================================================
\section{Discussion}\label{sec:discussion}
% 必须包含三块：(1) 方法与理论的局限（诚实、具体）；(2) 可扩展方向；(3) 实践建议。
% 禁止空话；每条局限要能对应一个可做的研究问题。

% =====================================================================
\bibliographystyle{chicago}   % JASA 风格；投其他刊按其要求切换
\bibliography{refs}

% =====================================================================
\appendix
% =====================================================================
% 附录：证明与补充理论（与 main.tex 配合编译）
% 统计学论文附录的标准结构：
%   A. 证明（按定理顺序，全部给出，不得引用不存在的中间引理）
%   B. 附加理论结果（备用引理 auxiliary lemmas，先证引理再证定理）
%   C. 额外模拟结果 / 附加真实数据分析
%   D. AI 辅助使用披露（本项目的特色，见 D 节模板）
% =====================================================================

\section{Proofs}\label{app:proof}
% 写证明前先跑 skills/stat-theory-check：
% 1) 每个引用的引理必须在此前或附录中已证；
% 2) 每个“显然/易得”必须展开或替换为可查引理；
% 3) 极限交换、求和交换必须有控制条件。

\subsection{Proof of Theorem \ref{thm:main}}\label{app:proof-thm1}
\begin{pf}
TODO：证明正文。典型结构：
Step 1（分解）：$\norm{\betaHat-\beta_0}$ 的基本不等式分解；
Step 2（控制逼近误差）：利用损失函数性质与 Assumption \ref{asm:1}；
Step 3（控制正则项误差）：KKT 条件 + 事件 $\mathcal{E}$ 的概率控制；
Step 4（合并）：在 $\lambda \asymp \sqrt{\log p/n}$ 下得到式 \eqref{eq:rate}。
（引用主文假设时用 \cref{asm:1}，须与主文实际存在的假设编号一致。）
\end{pf}

\section{Auxiliary lemmas}\label{app:lemmas}
% 所有被引用而未在正文出现的中间结果，按引理 A.1, A.2 编号并在证明前全部证完。

\section{Additional simulation results}\label{app:simsim}
% 补充场景、敏感性分析（如 λ 的取值网格、种子敏感性）。

\section{AI-assistance disclosure}\label{app:ai-disclosure}
% —— AI 辅助使用披露（按所在期刊政策调整；默认遵循 COPE 与各出版商 AI 使用规范）——
% 模板（逐条如实填写；未使用的项删除）：
%
% 1. 文献检索：使用 [工具名] 辅助检索与筛选文献，所有引用经人工逐条核对原文。
% 2. 代码：模拟与真实数据分析代码由作者编写，AI 用于代码解释与单元测试生成，
%    所有数值结果均由作者本地运行并复核（随机种子与环境见复现包）。
% 3. 理论推导：所有定理与证明由作者提出与验证；AI 用于检查证明中引理引用的一致性，
%    不作为正确性保证。
% 4. 语言润色：使用 [工具名] 对英文表达进行润色，不改变数学内容与结论。
% 5. AI 未被用于：数据伪造、结果生成、未经核验的引用插入。
%
% 提醒：多数期刊要求 AI 不能列为作者；使用披露需符合目标期刊的最新政策。
\end{document}
   % 全部证明、额外理论、补充图表（期刊正文页数限制内放不下的）

\newpage
% =====================================================================
% Supplement（补充材料，独立编译版）：用于期刊主线之外的
% 附加证明、附加模拟、附加数据集分析
% 若目标期刊不区分 appendix/supplement，可删除本文件并在 main.tex 中合并
% =====================================================================
\documentclass[11pt]{article}
\usepackage{amsmath,amssymb,amsthm}
\usepackage{booktabs,graphicx}
\usepackage[margin=1in]{geometry}
\usepackage[colorlinks=true,linkcolor=blue,citecolor=blue]{hyperref}
\usepackage{natbib}

% 与 main.tex 保持一致的定理计数（补充材料中定理编号前缀 S）
\newtheorem{sthm}{Theorem}
\renewcommand{\thesthm}{S\arabic{sthm}}
\newtheorem{slem}[sthm]{Lemma}
\renewcommand{\theslem}{S\arabic{sthm}}

\title{Supplement to ``TODO: 主文标题''}
\author{作者}
\date{}

\begin{document}
\maketitle

\tableofcontents

\section{Additional theoretical results}
% S1 定理及其证明；主文定理依赖的引理若篇幅过长，亦可放此。

\section{Additional simulation studies}
% 补充场景表：建议主文放代表性 2-3 场景，其余全量放补充材料。
% 每张表需在复现包中有对应脚本（sim/文件夹，随机种子固定）。

\section{Additional real data analysis}
% 第二个数据集的全部分析（统计论文顶刊常要求至少两个真实数据集或一个数据集+领域佐证）。

\section{Reproducibility statement}
% 复现包结构声明（TIER 4.0 风格）：
% data/      原始数据（或获取说明与哈希校验）
% code/      全部脚本，随机种子固定，requirements.txt/environment.yml 锁定环境
% output/    一键重新生成的全部图表
% 运行方式、预计时长、软硬件配置逐条列出。

\bibliographystyle{chicago}
\bibliography{refs}
\end{document}
 % 若期刊分主线+补充材料，则本文件独立编译为 supplement

\end{document}
